% ============================================================
%  ARTEMIS 课程实验报告正文
%  课程：3D 计算机图形学
%  题目：加速光线追踪与增强步进法在点云集成群聚算法中的应用
% ============================================================

\section{内容与要求}

\subsection{项目概述}

本项目名为 \textbf{ARTEMIS}（\textit{Accelerated Ray Tracing and Enhanced Marching
for Integrated flocking with point cloudS}），旨在构建一条从"文字或图像语义输入"
到"动态三维实体实时展示"的完整技术链路。系统将 AI 生成式点云、GPU 端群聚粒子模拟
与高性能光线渲染三条技术主线深度融合，最终实现成千上万个具备自主行为的粒子在屏幕上
动态"坍缩"成由人工智能生成的语义形状，并以物理正确的光照模型进行实时渲染。

项目整体可分解为以下四个子内容模块，各模块的功能要求如下。

% ------------------------------------------------------------
\subsection{子内容一：AI 驱动的三维点云生成管线}

\textbf{功能描述：}将人类的抽象语言或图像意向转化为三维离散坐标，作为粒子群落的
"终态蓝图"。

\textbf{所需功能：}
\begin{itemize}
  \item \textbf{文本到图像}：支持通过大语言模型（LLM）进行 Prompt Engineering，
        将自然语言描述转化为高质量参考图像（对接 Stable Diffusion / DALL-E 等模型）。
  \item \textbf{图像预处理}：对输入图像进行背景去除（\texttt{backgroundremover}），
        提取前景主体，提升点云生成质量。
  \item \textbf{点云扩散生成}：调用 OpenAI Point-E 模型（1B 参数量），
        以图像或文本为条件，通过扩散过程生成带有颜色信息的三维点云，
        输出格式为包含 $(x, y, z, r, g, b)$ 等特征的 \texttt{.npy} 文件，
        点数量在 1024 至 4096 之间。
  \item \textbf{点云筛选与后处理}：支持在多次生成结果中自动选取质量最优的点云
        （\texttt{automate\_selection.ipynb}），并提供点云到网格的转换接口
        （\texttt{pointcloud2mesh.ipynb}）以供可视化验证。
  \item \textbf{西安交通大学特色资产生成}：针对校徽、腾飞塑像、人工智能学院徽标等
        具有校园特色的目标物，探索适合 Point-E 的图像输入策略，
        生成可辨识的定制化点云资产。
\end{itemize}

% ------------------------------------------------------------
\subsection{子内容二：GPU 端 Boids 群聚粒子模拟}

\textbf{功能描述：}在 GPU 上以高度并行的方式驱动数千个粒子执行 Boids 群聚行为，
并使粒子最终向目标点云槽位汇聚。

\textbf{所需功能：}
\begin{itemize}
  \item \textbf{Reynolds 三定律}：在 OpenCL 内核中实现经典 Boids 的三大规则——
        分离（Separation）、凝聚（Cohesion）、对齐（Alignment），
        各规则具有独立的感知半径与权重系数。
  \item \textbf{Arrival 抵达规则}：为每个粒子分配目标点云中的一个槽位，
        通过速度缩放使粒子在接近目标时平滑减速，实现形态"坍缩"效果。
  \item \textbf{包围球约束}：防止粒子逃逸出场景范围，通过预测碰撞时刻
        施加切向偏转力。
  \item \textbf{数值积分}：采用带速度和加速度上限的半隐式 Verlet 积分，
        保证粒子运动稳定。
  \item \textbf{颜色插值}：粒子颜色随每帧向目标颜色线性插值过渡，
        实现从白色自由粒子到彩色点云形态的视觉渐变。
  \item \textbf{实时交互}：支持键盘切换目标点云模型、暂停/继续模拟、
        鼠标旋转与滚轮缩放视角。
\end{itemize}

% ------------------------------------------------------------
\subsection{子内容三：BVH 加速结构与 Morton 编码}

\textbf{功能描述：}在每帧渲染前对粒子进行空间划分，大幅减少光线求交的无效计算量。

\textbf{所需功能：}
\begin{itemize}
  \item \textbf{BVH 构建}：每帧在 CPU 侧对当前粒子位置递归构建层次包围盒（BVH），
        按最长轴排序分割，将 $N$ 个粒子组织为若干个大小为 \texttt{WORK\_GROUP\_SIZE}
        的叶节点，每个叶节点对应一个轴对齐包围盒（AABB）。
  \item \textbf{Morton 编码（备选加速路径）}：将粒子归一化坐标映射为
        30 位 Z-Order Morton 编码，使空间相邻粒子在显存中趋于物理连续，
        提升 GPU L1/L2 缓存命中率。
  \item \textbf{排列索引上传}：将 BVH 排序后的置换数组
        （\texttt{permutation}）与包围盒数组（\texttt{bboxes}）
        上传至 GPU 显存供渲染内核使用。
  \item \textbf{光线-AABB 快速剔除}：渲染内核先以位掩码方式标记光线与各叶节点
        包围盒的相交情况，仅对命中的叶节点内的粒子执行精确求交，跳过大量无关计算。
\end{itemize}

% ------------------------------------------------------------
\subsection{子内容四：实时光线渲染（Ray Tracing 与 Ray Marching）}

\textbf{功能描述：}对粒子群落进行物理正确的光照渲染，支持两种渲染模式，
通过 OpenCL–OpenGL 零拷贝互操作实现实时帧率输出。

\textbf{所需功能：}
\begin{itemize}
  \item \textbf{Ray Tracing 模式}：对每个粒子进行精确的球体光线求交，
        支持区域光源软阴影采样（\texttt{LIGHT\_SAMPLES} 次蒙特卡洛采样）
        和全局环境光照，分辨率最高支持 $3840 \times 2880$，粒子数量最高 4096。
  \item \textbf{Ray Marching 模式}：基于有向距离场（SDF）的步进渲染，
        通过 \texttt{smooth\_min} 平滑并集函数使相邻粒子表面自动融合为
        有机流体状曲面，迭代步数最多 40 步，最终通过梯度计算法线。
  \item \textbf{OpenCL–OpenGL Interop}：利用 \texttt{cl::ImageGL} 将
        OpenCL 计算结果直接写入 OpenGL 纹理，避免 CPU 内存与 GPU 显存之间的
        数据拷贝，消除 PCI-E 带宽瓶颈。
  \item \textbf{Work Group 局部内存优化}：在 GPU 内核中使用
        \texttt{local float3 localBuffer[WORK\_GROUP\_SIZE]} 将频繁访问的
        粒子数据缓存到 Work Group 共享内存，降低全局显存访问延迟。
  \item \textbf{跨平台支持}：针对 Linux（GLX）、Windows（WGL）、
        macOS（CGL ShareGroup）分别实现 OpenCL 上下文初始化，
        兼容主流操作系统与 GPU 平台。
\end{itemize}



% ============================================================
\section{总体方案}

\subsection{系统整体架构}

ARTEMIS 系统由两条相互解耦的子流水线组成，通过标准化的 \texttt{.npy} 点云文件
作为中间介质进行衔接，整体架构如图~\ref{fig:pipeline} 所示。

\begin{figure}[htbp]
  \centering
  \fbox{\rule{0pt}{6cm}\rule{12cm}{0pt}}
  \caption{ARTEMIS 系统整体流水线架构图（待替换为 figures/3d/pipeline.pdf）}
  \label{fig:pipeline}
\end{figure}

\begin{itemize}
  \item \textbf{生成侧（Python 管线）}：运行于配备 NVIDIA GeForce RTX 4090 D 显卡的
        Ubuntu Docker 服务器环境，负责从文本或图像出发，通过深度学习模型推理生成
        三维点云文件，输出至共享存储目录。
  \item \textbf{渲染侧（C++/OpenCL 管线）}：运行于搭载 Apple M4 Pro 芯片的
        MacBook 本地环境，读取点云文件后在 GPU 上实时执行 Boids 模拟与光线渲染，
        并将结果输出至屏幕。
\end{itemize}

\subsection{生成侧：AI 点云生成管线}

生成管线的完整流程为：
\[
  \text{自然语言 Prompt}
  \xrightarrow{\text{LLM}} \text{图像描述}
  \xrightarrow{\text{DALL-E/SD}} \text{RGB 图像}
  \xrightarrow{\text{backgroundremover}} \text{前景图像}
  \xrightarrow{\text{Point-E (1B)}} \texttt{.npy}
\]

Point-E 是 OpenAI 提出的基于扩散过程的三维点云生成模型，其核心思路是
先由图像条件生成低分辨率点云，再通过上采样网络细化。
本项目部署的是参数量为 \textbf{1B} 的 \texttt{image2pointcloud} 版本，
单次推理在 RTX 4090 D 上运行约 \textbf{【待填写】} 秒，
输出包含 $N=4096$ 个点的带色点云，每个点携带
$(x, y, z, r, g, b)$ 六维特征，以 \texttt{float32} 格式存储为 \texttt{.npy} 文件。

对于西安交通大学特色资产（校徽、腾飞塑像、人工智能学院徽标等），
针对 Point-E 对平面图案深度感知不足的问题，我们采用的优化策略为：
优先选取具有\textbf{清晰三维轮廓}的参考图像，
并在必要时借助 Stable Diffusion 对图像进行三维感知风格化处理后再输入 Point-E，
以提升生成点云的空间分布质量与可辨识度。

\subsection{渲染侧：C++/OpenCL 渲染管线}

渲染侧的帧循环结构如图~\ref{fig:frame_loop} 所示，
每一帧依次执行模拟（Simulate）、处理（Process）、渲染（Render）三个阶段。

\begin{figure}[htbp]
  \centering
  \fbox{\rule{0pt}{5cm}\rule{12cm}{0pt}}
  \caption{渲染侧帧循环结构示意图（待替换为 figures/3d/frame\_loop.pdf）}
  \label{fig:frame_loop}
\end{figure}

\subsubsection{阶段一：Boids 模拟（GPU 并行）}

\texttt{simulateTimeStep} 函数依次调用两个 OpenCL 内核：

\paragraph{calculate.cl — 加速度计算}
该内核在 GPU 上并行为每个粒子 $i$ 计算合力加速度，各分量定义如下。

\textbf{分离力（Separation）}：排斥感知半径 $r_s=0.1$ 内的相邻粒子，
\begin{equation}
  \mathbf{a}_{\mathrm{sep}} = K_s \cdot
  \widehat{\sum_{j:\,d_{ij}<r_s}
  \frac{\mathbf{p}_i - \mathbf{p}_j}{d_{ij}^2}}, \quad K_s = 7
\end{equation}

\textbf{凝聚力（Cohesion）}：驱动粒子向感知半径 $r_c=0.15$ 内邻居质心靠拢，
\begin{equation}
  \mathbf{a}_{\mathrm{coh}} = K_c \cdot
  \widehat{v_{\max}\widehat{\bar{\mathbf{p}}_{\mathrm{nbr}}-\mathbf{p}_i}
  - \mathbf{v}_i}, \quad K_c = 7.5
\end{equation}

\textbf{对齐力（Alignment）}：使粒子速度向感知半径 $r_a=0.08$ 内邻居平均速度对齐，
\begin{equation}
  \mathbf{a}_{\mathrm{aln}} = K_a \cdot
  \widehat{v_{\max}\widehat{\bar{\mathbf{v}}_{\mathrm{nbr}}} - \mathbf{v}_i},
  \quad K_a = 4
\end{equation}

\textbf{包围球约束（Containment）}：预测粒子碰撞包围球（$R=1$）的时刻 $t^*$，
施加切向偏转力。令 $a=\|\mathbf{v}\|^2$，$b=2\mathbf{p}\cdot\mathbf{v}$，
$c=\|\mathbf{p}\|^2-R^2$，则
\begin{equation}
  t^* = \frac{\sqrt{b^2-4ac}-b}{2a}, \qquad
  \mathbf{a}_{\mathrm{con}} = K_{\mathrm{con}} \cdot
  \widehat{\mathbf{n}-(\mathbf{n}\cdot\hat{\mathbf{v}})\hat{\mathbf{v}}},
  \quad K_{\mathrm{con}}=20
\end{equation}

\textbf{抵达力（Arrival）}：每个粒子被分配目标点云中唯一槽位 $\mathbf{g}_i$，
通过速度缩放平滑减速，
\begin{equation}
  \mathbf{v}^*_{\mathrm{arr}} =
  \min\!\left(v_{\max},\,
  \frac{\|\mathbf{g}_i-\mathbf{p}_i\|}{r_{\mathrm{arr}}}\right)
  \widehat{\mathbf{g}_i-\mathbf{p}_i}, \quad
  \mathbf{a}_{\mathrm{arr}} = K_{\mathrm{arr}} \cdot
  \widehat{\mathbf{v}^*_{\mathrm{arr}}-\mathbf{v}_i},
  \quad K_{\mathrm{arr}}=70,\; r_{\mathrm{arr}}=0.05
\end{equation}

\paragraph{integrate.cl — 数值积分}
采用带速度上限的半隐式 Verlet 积分更新粒子状态：
\begin{equation}
  \mathbf{p}_i \leftarrow \mathbf{p}_i + \Delta t\,\mathbf{v}_i
  + \tfrac{\Delta t^2}{2}\,\mathbf{a}_i, \qquad
  \mathbf{v}_i \leftarrow \mathbf{v}_i + \Delta t\,\mathbf{a}_i,
  \quad \|\mathbf{v}_i\| \leq v_{\max}=0.5
\end{equation}
粒子颜色同步按
$\mathbf{c}_i \leftarrow (1-\Delta t)\mathbf{c}_i + \Delta t\,\mathbf{c}^{\mathrm{target}}_i$
逐帧线性插值。

\subsubsection{阶段二：BVH 构建与数据上传（CPU）}

\texttt{processTimeStep} 函数首先从 GPU 读回当前粒子坐标，在 CPU 上递归构建
BVH：每次选取粒子分布范围最长的坐标轴进行排序分割，
直至每个叶节点恰好包含 \texttt{WORK\_GROUP\_SIZE}（=64）个粒子。
核心伪代码如下：

\begin{verbatim}
void bvh(spheres, permutation, l, r):
    if r - l <= WORK_GROUP_SIZE: return
    axis = argmax(max_x-min_x, max_y-min_y, max_z-min_z)
    sort permutation[l..r) by spheres[axis]
    mid = (l + r) / 2
    bvh(spheres, permutation, l, mid)
    bvh(spheres, permutation, mid, r)
\end{verbatim}

BVH 完成后，计算每个叶节点的轴对齐包围盒（AABB），
将置换数组 \texttt{permutation} 与包围盒数组 \texttt{bboxes}
通过 \texttt{enqueueWriteBuffer} 上传至 GPU 供渲染内核使用。

\subsubsection{阶段三-A：Ray Tracing 模式}

对每条相机光线依次执行：
\begin{enumerate}
  \item 以 64 位掩码标记与各叶节点 AABB 相交的情况（BVH 剔除）；
  \item 将命中叶节点内的粒子坐标加载至局部共享内存 \texttt{localBuffer}；
  \item 对局部内存中每个球体执行精确光线-球求交，
        判别式 $\Delta=b^2-4ac\geq0$ 时取最近交点 $t_{\min}$；
  \item 在命中点向区域光源发射 \texttt{LIGHT\_SAMPLES} 条阴影光线，
        按蒙特卡洛面光源采样公式计算辐亮度贡献：
        \begin{equation}
          L_o = \frac{1}{N_s}\sum_{k=1}^{N_s}
          \frac{f_r \cdot L_e \cdot \cos\theta_i}{p(\omega_k)},
          \quad p(\omega_k) = \frac{r^2}{A\,|\cos\theta_e|}
        \end{equation}
\end{enumerate}
Ray Tracing 模式下分辨率最高支持 $3840\times2880$，
粒子数 $N=4096$，光源采样数 \texttt{LIGHT\_SAMPLES}$=8$。

\subsubsection{阶段三-B：Ray Marching 模式}

基于 SDF 的球追踪（Sphere Tracing）算法：
\begin{enumerate}
  \item 对当前步进点 $\mathbf{h}$ 计算带平滑并集的 SDF 值：
        \begin{equation}
          d(\mathbf{h}) =
          \mathop{\mathrm{smooth\_min}}_{i}
          \bigl(\|\mathbf{h}-\mathbf{p}_i\|-r_{\mathrm{sphere}}\bigr)
        \end{equation}
        \begin{equation}
          \mathrm{smooth\_min}(a,b)=\min(a,b)-\frac{h^3 k}{6},\quad
          h=\max\!\left(\frac{k-|a-b|}{k},0\right),\quad k=0.06
        \end{equation}
  \item 步进 $t\leftarrow t+d(\mathbf{h})$，最多迭代 40 步；
  \item 命中（$d<0.001$）后用中心差分估算法线：
        \begin{equation}
          \mathbf{n} \approx \nabla d(\mathbf{h}) \approx
          \frac{1}{2\varepsilon}
          \bigl(d(x+\varepsilon)-d(x-\varepsilon),\;
                d(y+\varepsilon)-d(y-\varepsilon),\;
                d(z+\varepsilon)-d(z-\varepsilon)\bigr),
          \quad \varepsilon=0.001
        \end{equation}
  \item 以法线结合面光源采样计算漫反射辐亮度。
\end{enumerate}

Ray Marching 模式下粒子数 $N=1024$，球体半径 $r_{\mathrm{sphere}}=0.015$，
分辨率 $800\times600$。\texttt{smooth\_min} 使相邻粒子在空间中自动融合，
视觉上呈现出磁流体般的有机连续曲面，是 ARTEMIS 区别于普通粒子系统的核心视觉特征。

\subsubsection{OpenCL–OpenGL 零拷贝互操作}

渲染结果通过 \texttt{cl::ImageGL} 直接写入 OpenGL 纹理对象，
无需经过 CPU 内存中转：
\texttt{glFinish()} $\to$ \texttt{enqueueAcquireGLObjects()} $\to$
渲染内核写纹理 $\to$ \texttt{enqueueReleaseGLObjects()} $\to$
OpenGL 全屏四边形输出至窗口。
此机制消除了 PCI-E 带宽瓶颈，是实现千粒子实时渲染的重要工程基础。

\subsection{硬件与软件环境总结}

\begin{table}[htbp]
  \centering
  \caption{实验硬件与软件环境汇总}
  \label{tab:env}
  \begin{tabular}{p{2.2cm}p{5cm}p{6.5cm}}
    \hline
    \textbf{用途} & \textbf{硬件平台} & \textbf{软件栈} \\
    \hline
    点云模型推理 &
      NVIDIA GeForce RTX 4090 D（49 GB），
      8 核 CPU，Ubuntu 22.04 Docker &
      Python 3.9，PyTorch 2.5.1+cu121，Point-E 1B，
      backgroundremover，JupyterLab \\
    \hline
    实时仿真渲染 &
      Apple MacBook，M4 Pro 芯片 &
      macOS，Clang/C++17，OpenCL 1.2，
      OpenGL 4.1，GLFW 3.3，CMake 3.16+ \\
    \hline
  \end{tabular}
\end{table}


% ============================================================
\section{实验结果}

\subsection{实验环境与测试配置}

本实验在两台设备上分阶段完成。点云生成阶段在配备单卡
NVIDIA GeForce RTX 4090 D（显存 49 GB）的 Ubuntu 22.04 Docker
服务器上运行，使用 CUDA 12.1 与 PyTorch 2.5.1 驱动 Point-E 1B 模型；
实时渲染阶段在搭载 Apple M4 Pro 芯片的 MacBook 上完成，
通过 OpenCL 1.2 + OpenGL 4.1 实现 GPU 端 Boids 仿真与光线渲染。
系统关键参数汇总如表~\ref{tab:params} 所示。

\begin{table}[htbp]
  \centering
  \caption{系统关键参数汇总}
  \label{tab:params}
  \begin{tabular}{lll}
    \hline
    \textbf{参数} & \textbf{Ray Tracing 模式} & \textbf{Ray Marching 模式} \\
    \hline
    粒子数量 $N$ & 4096 & 1024 \\
    渲染分辨率 & $3840\times2880$ & $800\times600$ \\
    球体半径 & 0.01 & 0.015 \\
    Work Group 大小 & $8\times8=64$ & $8\times8=64$ \\
    每像素光线数 & 1 & 1 \\
    光源采样数 & 8 & 1 \\
    smooth\_min 系数 $k$ & — & 0.06 \\
    BVH 叶节点大小 & 64 & 64 \\
    平均帧率（M4 Pro） & 【待填写】 FPS & 【待填写】 FPS \\
    单帧 BVH 构建耗时 & 【待填写】 $\mu$s & 【待填写】 $\mu$s \\
    \hline
  \end{tabular}
\end{table}

\subsection{AI 点云生成结果}

图~\ref{fig:pointcloud_natural} 展示了以自然生物与日常物体为目标的
Point-E 生成结果，图中左侧为原始参考图像，右侧为生成的三维点云着色渲染图。

\begin{figure}[htbp]
  \centering
  % 建议采用 2 列 pair 排列：原图 | 点云渲染图
  \fbox{\rule{0pt}{5cm}\rule{14cm}{0pt}}
  \caption{自然物体点云生成结果示例（左：原始输入图像；右：Point-E 生成点云渲染）。
           图片目录：\texttt{figures/3d/pointcloud\_natural.png}（待替换）}
  \label{fig:pointcloud_natural}
\end{figure}

图~\ref{fig:pointcloud_xjtu} 展示了以西安交通大学特色元素为目标的定制化点云资产，
包括腾飞塑像、人工智能学院徽标等，体现了系统对校园文化特色的个性化适配能力。

\begin{figure}[htbp]
  \centering
  \fbox{\rule{0pt}{5cm}\rule{14cm}{0pt}}
  \caption{西安交通大学特色资产点云生成结果（左：参考图像；右：生成点云渲染）。
           图片目录：\texttt{figures/3d/pointcloud\_xjtu.png}（待替换）}
  \label{fig:pointcloud_xjtu}
\end{figure}

Point-E 1B 模型单次推理耗时约 \textbf{【待填写】} 秒，
生成的每份点云包含 4096 个带色离散点，
对于具有清晰三维轮廓的输入图像（如生物、雕塑等），
点云空间分布质量良好，颜色映射准确；
对于高度对称的平面图案（如校徽），
输出点云深度分布较为扁平，
具体改进方案详见第~\ref{sec:reflection} 节的分析。

\subsection{Boids 群聚动态演化过程}

图~\ref{fig:boids_process} 展示了粒子群从初始随机分布状态到最终收敛至目标
点云形态的动态演化过程，共截取四个关键时刻。

\begin{figure}[htbp]
  \centering
  \fbox{\rule{0pt}{6cm}\rule{14cm}{0pt}}
  \caption{Boids 群聚动态演化过程（从左至右依次为：初始随机态、开始聚集、
           中间过渡、最终收敛）。图片目录：\texttt{figures/3d/boids\_process.png}（待替换）}
  \label{fig:boids_process}
\end{figure}

在实验中观察到，Arrival 力系数 $K_{\mathrm{arr}}=70$ 相较于其余 Boids 力
的量级明显偏大，这使得粒子在启动目标追踪后能够快速压制随机游走行为，
在约 \textbf{【待填写】} 秒内完成从散乱态到聚集态的视觉收敛。
Separation 与 Cohesion 的协同作用则确保粒子在收敛过程中不会堆叠重叠，
维持一定的空间密度分布。

\subsection{Ray Tracing 渲染效果}

图~\ref{fig:raytrace_demo} 展示了 Ray Tracing 模式下多个点云目标的渲染效果。
区域光源的软阴影与环境光照营造出真实感较强的三维立体效果，
粒子的颜色信息由 Point-E 生成的 RGB 特征直接映射，
视觉上呈现出"彩色点云立体浮雕"的效果。

\begin{figure}[htbp]
  \centering
  \fbox{\rule{0pt}{6cm}\rule{14cm}{0pt}}
  \caption{Ray Tracing 模式渲染效果 demo（分辨率 $3840\times2880$，$N=4096$）。
           图片目录：\texttt{figures/3d/raytrace\_demo.png}（待替换）}
  \label{fig:raytrace_demo}
\end{figure}

\subsection{Ray Marching 渲染效果}

图~\ref{fig:raymarch_demo} 展示了 Ray Marching 模式下的渲染结果。
相较于 Ray Tracing 的离散球体效果，Ray Marching 通过 \texttt{smooth\_min}
平滑并集将相邻粒子的 SDF 表面自然融合，视觉上粒子群落呈现出
流体状的有机连续曲面，更具视觉冲击力。

\begin{figure}[htbp]
  \centering
  \fbox{\rule{0pt}{6cm}\rule{14cm}{0pt}}
  \caption{Ray Marching 模式渲染效果 demo（分辨率 $800\times600$，$N=1024$）。
           图片目录：\texttt{figures/3d/raymarch\_demo.png}（待替换）}
  \label{fig:raymarch_demo}
\end{figure}

图~\ref{fig:comparison} 将两种渲染模式进行对比，
直观呈现 \texttt{smooth\_min} 带来的表面融合效果与 Ray Tracing 离散球体效果的差异。

\begin{figure}[htbp]
  \centering
  \fbox{\rule{0pt}{5cm}\rule{14cm}{0pt}}
  \caption{渲染模式对比（左：Ray Tracing 离散球体；右：Ray Marching 平滑融合）。
           图片目录：\texttt{figures/3d/comparison.png}（待替换）}
  \label{fig:comparison}
\end{figure}

\subsection{性能分析}

表~\ref{tab:perf} 汇总了各阶段的时间开销统计（基于代码中
\texttt{PROFILE} 宏的计时结果）。

\begin{table}[htbp]
  \centering
  \caption{各阶段平均耗时统计（M4 Pro，单位 $\mu$s）}
  \label{tab:perf}
  \begin{tabular}{lcc}
    \hline
    \textbf{阶段} & \textbf{Ray Tracing 模式} & \textbf{Ray Marching 模式} \\
    \hline
    Boids 模拟（simulate）   & 【待填写】 & 【待填写】 \\
    BVH 构建（CPU）          & 【待填写】 & 【待填写】 \\
    渲染内核（process）      & 【待填写】 & 【待填写】 \\
    OpenGL 输出（render）    & 【待填写】 & 【待填写】 \\
    总帧耗时 / 帧率           & 【待填写】 $\mu$s / 【待填写】 FPS
                             & 【待填写】 $\mu$s / 【待填写】 FPS \\
    \hline
  \end{tabular}
\end{table}

从性能结构来看，BVH 构建当前在 CPU 侧执行，
是整条管线中潜在的主要瓶颈之一：在粒子数 $N=4096$ 时，
该步骤需要每帧对粒子坐标排序，复杂度为 $O(N\log N)$，
并伴随 \texttt{enqueueReadBuffer} 的 GPU-CPU 数据回读开销。
而渲染内核在 GPU 上借助 BVH 剔除与 Work Group 局部内存优化，
实际每线程的球体求交次数大幅低于暴力 $O(N)$ 遍历，
是支撑实时帧率的核心加速手段。


% ============================================================
\section{心得体会}
\label{sec:reflection}

\subsection{整体架构设计的思考}

接手这个项目的架构设计工作，说实话一开始并没有完全想清楚整条技术路线应该怎么串联。
最早的想法很简单——把 Point-E 生成的点云直接用来做 Boids 的目标点，然后渲染出来。
但随着对 OpenCL 和 OpenGL 互操作机制的深入理解，我意识到
"如何让 GPU 上的计算结果不经过 CPU 中转直接成为下一阶段的输入"
才是决定系统性能上限的核心问题。
\texttt{cl::ImageGL} 零拷贝互操作这个设计决策是整个架构里我最满意的一个地方，
它把"Boids 模拟输出 $\to$ 光线渲染输入 $\to$ OpenGL 显示"这三段天然地焊接在了一起，
消除了本来会出现在每一帧里的大量冗余内存拷贝。

另一个比较重要的架构决策是将系统切分为"生成侧（Python/服务器）"和
"渲染侧（C++/本地）"两条完全解耦的流水线，通过 \texttt{.npy} 文件作为
唯一的接口契约。这样做的好处是两条线可以独立开发和调试：
服务器上跑模型推理不需要关心渲染细节，
本机渲染也不需要等待每次推理完成才能测试。
事后来看，这个解耦设计替我们节省了大量来回联调的时间。

\subsection{点云模型推理与渲染调优的经验}

在点云生成这一侧，踩坑最多的地方是\textbf{输入图像质量对 Point-E 输出的影响}。
表~\ref{tab:trial_error} 总结了主要的试错经验。

\begin{table}[htbp]
  \centering
  \caption{点云生成主要 Trial \& Error 汇总}
  \label{tab:trial_error}
  \begin{tabular}{p{3cm}p{4.5cm}p{6cm}}
    \hline
    \textbf{问题现象} & \textbf{根本原因} & \textbf{解决策略} \\
    \hline
    校徽点云几乎全在一个平面上 &
      Point-E 对纯 2D 图案深度感知极差，
      缺乏三维线索 &
      改用具有立体感的实物照片或
      借助 Stable Diffusion 做三维感知风格化处理后再输入 \\
    \hline
    全景图点云细节糊成一团 &
      点云数量有限（1024～4096），
      信息密度过高导致细节丢失 &
      裁剪聚焦主体，去除背景，
      保证前景轮廓清晰 \\
    \hline
    Flocking 收敛后形态不可辨识 &
      粒子群需要清晰的三维轮廓才能
      聚集出可辨识形态 &
      优先选取三维轮廓感强的目标
      （雕塑、生物侧面等），
      放弃对称平面徽标 \\
    \hline
    颜色混杂，标志性色彩丢失 &
      颜色区域太碎，点云颜色映射
      失去辨识度 &
      选取颜色区域集中、对比度高的
      图像，或在后处理中对颜色聚类 \\
    \hline
  \end{tabular}
\end{table}

最开始尝试直接用西安交通大学校徽图案生成点云时，
得到的结果几乎是一张"薄饼"——所有点云都挤在一个平面上，
Boids 聚集后的形态从任何角度看都像是贴在屏幕上的一个平面图案，
完全没有三维感。
这让我意识到 Point-E 本质上是在做"从 2D 推断 3D"，
如果输入图像本身就没有深度信息，模型自然也无从猜测。
后来换用腾飞塑像的侧面实物照片后，效果明显改善，
生成的点云有了真实的三维分布，Boids 聚集时的视觉效果也好了很多。

在渲染调优这侧，\textbf{Boids 参数调整}是另一个反复折腾的过程。
Arrival 力系数 $K_{\mathrm{arr}}=70$ 是经过多次实验最终确定的值——
值太小时粒子在边游荡边向目标靠拢，形态收敛极慢，演示效果差；
值太大时粒子几乎直线冲向目标，失去了群聚的有机感，像一堆子弹而非鸟群。
70 这个值在"快速收敛"和"保留 Boids 有机运动质感"之间取得了相对平衡的 trade-off。

Separation 半径 $r_s=0.1$ 和 Cohesion 半径 $r_c=0.15$ 的选取同样经历了反复迭代。
早期两个半径相差不大时，粒子会形成极度紧密的团块，
Ray Marching 的 \texttt{smooth\_min} 将它们渲染成一个巨大的"肉球"，
失去了点云的形态特征。
适当拉开两个半径的差距后，粒子群在保持整体凝聚性的同时维持了一定的颗粒感，
形态可辨识度大幅提升。

\subsection{两种渲染模式的 Trade-off 分析}

Ray Tracing 与 Ray Marching 两种模式各有优劣，
在实验过程中积累了一些直观的认识：

\begin{itemize}
  \item \textbf{Ray Tracing}：每个粒子是独立的精确球体，颜色信息保留完整，
        在高分辨率（$3840\times2880$）下视觉效果细腻，
        适合展示点云的颜色分布与空间结构细节。
        代价是粒子数量受限（$N=4096$），
        且粒子之间没有表面融合，整体感稍弱。
  \item \textbf{Ray Marching}：通过 \texttt{smooth\_min} 实现表面融合，
        群落整体呈现出流体状有机曲面，视觉冲击力更强，
        特别适合展示 Boids 聚集的动态过程——粒子从散到聚的过程中，
        曲面像液体一样逐渐成形，非常震撼。
        代价是分辨率与粒子数量受限（$N=1024$，$800\times600$），
        且无法独立保留每个粒子的颜色信息。
\end{itemize}

在实际演示中，我们的策略是先以 Ray Marching 模式展示群聚收敛的震撼过程，
再切换至 Ray Tracing 模式展示点云的颜色细节，形成互补的演示节奏。

\subsection{对 GPU 并行编程的理解}

通过这次课程设计，对 GPU 并行编程有了远比课堂上更深的体感认识。
以前觉得"把循环并行化"就是 GPU 编程的全部，
但在实际实现 Boids 内核时才意识到，
\textbf{访存模式}才是决定 GPU 程序性能的第一要素：
同一个 Work Group 内的线程需要访问哪些数据？
这些数据是否在物理上连续？能否装进局部内存（Local Memory）？

\texttt{local float3 localBuffer[WORK\_GROUP\_SIZE]} 这段代码看起来只是一个
数组声明，背后却代表着一次对 GPU 存储层次的主动利用——
把每个 Work Group 需要反复访问的粒子坐标提前"拉"到片上共享内存，
避免每次都去高延迟的全局显存（Global Memory）取数据。
Morton 编码（Z-Order 排序）的思路也类似：
让空间上相邻的粒子在显存中也连续存放，从而提升 L1/L2 缓存的命中率。
这些优化手段在课本里看起来都是"术"，
但只有真正在代码里量化出 BVH 构建耗时、渲染内核耗时的对比数字之后，
才能切实感受到它们的价值。

\subsection{不足与展望}

在整个实验过程中，也清醒地认识到了若干未竟之处：

\begin{enumerate}
  \item \textbf{BVH 未能 GPU 化}：当前 BVH 构建在 CPU 侧执行，
        每帧需要 \texttt{enqueueReadBuffer} 将粒子坐标从 GPU 读回 CPU，
        再排序后重新上传，这一来回造成了不可忽视的延迟。
        若能将 BVH 构建完全移植到 OpenCL 内核（基于 Morton 编码的并行基数排序），
        预计可显著降低每帧的 process 阶段耗时。
  \item \textbf{Boids 计算复杂度}：当前 \texttt{calculate.cl} 对每个粒子遍历全部
        $N$ 个邻居，复杂度为 $O(N^2)$。随着粒子数增加，
        仿真负载将呈平方增长。引入网格空间哈希（Grid Spatial Hashing）
        可将检索复杂度降至接近 $O(N)$，是未来的主要优化方向。
  \item \textbf{macOS OpenCL 兼容性}：Apple 已在新版 macOS 上废弃 OpenCL，
        在 M4 Pro 上运行时会出现 API 飘黄警告，且无法利用 Metal 的最新硬件特性。
        未来若要在 Apple Silicon 上获得最优性能，应当考虑迁移至 Metal Compute 或
        WebGPU 框架。
  \item \textbf{点云生成质量}：Point-E 在处理平面图案和细节丰富的场景时
        生成质量有限，受限于模型本身的能力边界。
        若引入更新的三维生成模型（如 Shap-E、Zero123 等），
        有望在相同输入质量下获得更好的点云分布。
\end{enumerate}

总的来说，这次课程设计让我真正体会到了"从零到一"地将图形学理论落地为可运行系统的
完整过程——从最初在白板上画 pipeline 框图，到在服务器上跑通 Point-E 推理，
再到在本机上看着几千个粒子慢慢聚拢成一个三维形态——每一个阶段都有意外的困难，
也都有解决问题之后那种说不清楚的满足感。
光线追踪、SDF 渲染、GPU 并行编程这些在课堂上听起来各自独立的知识点，
在这个项目里被强行"焊"在了一起，反而让人觉得它们本来就应该是一体的。
